% Source: https://artofproblemsolving.com/community/c3h2827341_area_of_an_ellipse
% Stable identity: aops-topic-2827341-q01
The map $(x,y)\mapsto\left(\dfrac abx,\,y\right)$ sends the circle $x^2+y^2=b^2$ onto the ellipse, since $\dfrac{(ax/b)^2}{a^2}+\dfrac{y^2}{b^2}=\dfrac{x^2+y^2}{b^2}$.
\begin{asy}
usepackage("amsmath");
size(6cm);
defaultpen(fontsize(10pt));
import graph;draw(scale(2,1)*unitcircle);draw(unitcircle,dashed);
draw((-2.4,0)--(2.4,0),Arrow);draw((0,-1.4)--(0,1.4),Arrow);
label("$a$",(1.5,0),S);label("$b$",(0,0.5),W);label("$x$-scale by $a/b$",(1.3,0.85),NE);
draw((0.6,0.8)--(1.2,0.8),Arrow);
\end{asy}

Stretching in one direction by the factor $\dfrac ab$ multiplies every area by $\dfrac ab$: every horizontal strip of the region keeps its height and has its length multiplied by $\dfrac ab$ (Cavalieri's principle). So the area of the ellipse is $\dfrac ab\cdot\pi b^2=\boxed{\pi ab}$. $\blacksquare$
